% from ega0-1-fr.tex
\oldpage[0\textsubscript{I}]{53}

En effet, pour tout $\mathcal{O}_Y$-Module $\mathcal{L}$, on a un
homomorphisme canonique
\[
  f_*(\mathcal{F})\otimes_{\mathcal{O}_Y}\mathcal{L}
  \xrightarrow{1\otimes\rho}
  f_*(\mathcal{F})\otimes_{\mathcal{O}_Y}f_*(f^*(\mathcal{L}))
  \xrightarrow{\alpha}
  f_*\bigl(\mathcal{F}\otimes_{\mathcal{O}_X}f^*(\mathcal{L})\bigr),
\]
$\rho$ étant l'homomorphisme (4.4.3.2) et $\alpha$ l'homomorphisme (4.2.2.1).
Pour montrer que lorsque $\mathcal{L}$ est localement libre, cet homomorphisme
est bijectif, il suffit, la question étant locale sur $Y$, de considérer le
cas où $\mathcal{L}=\mathcal{O}_Y^n$~; en outre, $f_*$ et $f^*$ commutant aux
sommes directes finies, on peut supposer $n=1$, et dans ce cas la proposition
résulte aussitôt des définitions et de la relation
$f^*(\mathcal{O}_Y)=\mathcal{O}_X$.
\end{env}

\subsection{Faisceaux sur un espace annelé en anneaux locaux}
\label{subsection:0.5.5-fr}

\begin{env}[5.5.1]
\label{0.5.5.1-fr}
Nous dirons qu'un espace annelé $(X,\mathcal{O}_X)$ est un \emph{espace
annelé en anneaux locaux} si, pour tout $x\in X$, $\mathcal{O}_x$ est un
anneau local~; ces espaces annelés seront de loin les plus fréquents de ceux
qui seront considérés dans ce travail. Nous désignerons alors par
$\mathfrak{m}_x$, l'\emph{idéal maximal} de $\mathcal{O}_x$, par $k(x)$ le
\emph{corps résiduel} $\mathcal{O}_x/\mathfrak{m}_x$~; pour tout
$\mathcal{O}_X$-Module $\mathcal{F}$, tout ouvert $U$ de $X$, tout point
$x\in U$ et toute section $f\in\Gamma(U,\mathcal{F})$, nous désignerons par
$f(x)$ la \emph{classe} du germe
$f_x\in\mathcal{F}_x\pmod{\mathfrak{m}_x\mathcal{F}_x}$, et nous dirons que
c'est la \emph{valeur} de $f$ au point $x$. La relation $f(x)=0$ signifie
donc que $f_x\in\mathfrak{m}_x\mathcal{F}_x$~; lorsqu'elle est remplie, on dira
(par abus de langage) que $f$ \emph{s'annule en $x$}. On aura soin de ne pas
confondre cette relation avec $f_x=0$.
\end{env}

\begin{env}[5.5.2]
\label{0.5.5.2-fr}
Soient $X$ un espace annelé en anneaux locaux, $\mathcal{L}$ un
$\mathcal{O}_X$-Module inversible, $f$ une section de $\mathcal{L}$ au-dessus
de $X$. Il y a alors \emph{équivalence} entre les trois propriétés suivantes
en un point $x\in X$~:
\begin{enumerate}
  \item[$a)$] \emph{$f_x$ est un générateur de $\mathcal{L}_x$}~;
  \item[$b)$] \emph{$f_x\notin\mathfrak{m}_x\mathcal{L}_x$} (autrement dit,
    $f(x)\neq0$)~;
  \item[$c)$] \emph{Il existe une section $g$ de $\mathcal{L}^{-1}$ au-dessus
    d'un voisinage ouvert $V$ de $x$ telle que l'image canonique de $f\otimes
    g$ dans $\Gamma(V,\mathcal{O}_X)$ (5.4.3) soit la section unité}.
\end{enumerate}

En effet, la question étant locale, on peut se borner au cas où
$\mathcal{L}=\mathcal{O}_X$~; l'équivalence de $a)$ et $b)$ est alors
évidente, et il est clair que $c)$ entraîne $b)$. D'autre part, si
$f_x\notin\mathfrak{m}_x$, $f_x$ est inversible dans $\mathcal{O}_x$, soit
$f_xg_x=1_x$. Par définition des germes de sections, cela veut dire qu'il
existe un voisinage $V$ de $x$ et une section $g$ de $\mathcal{O}_X$ au-dessus
de $V$ telle que $fg=1$ dans $V$, d'où $c)$.

Il résulte aussitôt de la condition $c)$ que l'ensemble $X_f$ des $x$
vérifiant les conditions équivalentes $a)$, $b)$, $c)$ est \emph{ouvert} dans
$X$~; suivant la terminologie introduite dans (5.5.1), c'est l'ensemble des
$x$ où $f$ \emph{ne s'annule pas}.
\end{env}

\begin{env}[5.5.3]
\label{0.5.5.3-fr}
Sous les hypothèses de (5.5.2), soit $\mathcal{L}'$ un second
$\mathcal{O}_X$-Module inversible~; alors, si
$f\in\Gamma(X,\mathcal{L})$, $g\in\Gamma(X,\mathcal{L}')$, on a
\[
  X_f\cap X_g=X_{f\otimes g}.
\]
On se ramène en effet aussitôt au cas où
$\mathcal{L}=\mathcal{L}'=\mathcal{O}_X$ (la question étant locale)~; comme
$f\otimes g$ s'identifie alors canoniquement au produit $fg$, la proposition
est évidente.
\end{env}

